\documentclass[]{../../../tex/classes/styledArticle}
\usepackage{../../../tex/packages/hebrewSupport}

\usepackage{upgreek}
%\usepackage{txfonts}
\usepackage{countriesofeurope}
\input pdfmsym
%\newcommand \bs \blacksquare

\author{שחר פרץ}
\title{\textit{סקירה בלשון} (אבל מודפס)}
\begin{document}
	\maketitle
	
	\large 
	האקדמיה ללשון העברית היא מוסד מדע ותרבות הפועל מכוח החוק. יש לו מספר מטרות; לדוגמה, עיצוב, פיתוח והכוונת הלשון העברית (בן־יהודה, 2024). היא גם ממציאה מילים חדשות על־מנת לחדש את השפה, וחוקרת את הספרות העברית לדורותיה (בן־יהודה, 2025, וכן גדיש, 2019) במטרה לבנות קונקורנציה מקיפה לשפה העברית (גדיש, 2019). 
	
	האקדמיה ללשון העברית התפתחה רבות בעשורים האחרונים; היא חברה פרקים בתחום הדקדוק וכללי הכתב, פועלת להקמת מנווה האקדמיה החדש, וכן לאחרונה הקימה מרכז הדרכה והשתלמויות בעברית (בן־יהודה, 2024). בשנים האחרונות היא אף הרחיבה את פעילותה לכיוונים נוספים: פעילויות לציבור הרחב, פעילות ענפה במרשתת, ופניה לדור הצעיר באמצעות הרדתות החברתיות (לפי בן־יהודה, 2024, וגדיש, 2019). נבחין שגדיש (2019) מרחיבה שהאקדמיה הקימה בשנים האחרונות מערכת פניה יעדותית, אתר עשיר, ויישומונים (אפליקציות). ישנן גם מילים חדשות שהאקדמיה מנסה לשלב בשפה המדוברת (גדיש, 2019, בן־יהודה, 2024). 
	
	\textbf{ביבליוגרפיה: }
	\begin{itemize}
		\item בן־יהודה ט. (2024). ''\textit{בת שבעים וצופה על הבאות}``. האקדמיה ללשון העברית. 
		\item גדיש ר. (2019). ''\textit{השירות האמיתי של האקדמיה ללשון העברית, וכך צריך להיות}``. מאקו. 
	\end{itemize}
	
	
\end{document}